% chapters/12/rangkuman.tex
\bagian{Rangkuman}

\begin{rangkumanbox}
\begin{itemize}
  \item Mesin permainan mengambil alih siklus permainan, perhitungan fisika,
        deteksi tabrakan, penayangan, dan penyuntingan level. Kode Anda
        berpindah menjadi mengisi tahap yang disediakan mesin.
  \item Unity dipasang melalui dua program: \istilah{Unity Hub} untuk mengurus
        versi dan \istilah{project}, serta \istilah{Unity Editor} sebagai tempat
        bekerja. Versi \istilah{LTS} dipilih karena lebih stabil dan lebih
        banyak dipakai.
  \item Empat panel utama: \istilah{Hierarchy} memuat benda yang ada di dalam
        scene, \istilah{Scene} menampilkan susunan ruangnya, \istilah{Inspector}
        menampilkan komponen benda yang dipilih, dan \istilah{Project window}
        memuat seluruh berkas project.
  \item \istilah{GameObject} adalah wadah kosong yang menjadi berguna karena
        komponen yang menempel padanya. Komponen \istilah{Transform} wajib ada;
        komponen lain ditambahkan sesuai keperluan, termasuk skrip C\# Anda.
  \item Skrip C\# menjadi komponen yang menempel pada benda. Satu skrip dapat
        dipakai oleh banyak benda, dan masing-masing menjalankannya sendiri.
  \item Fungsi skrip dipanggil mesin pada waktu yang telah ditentukan:
        \code{Awake} dan \code{Start} sekali di awal, \code{Update} setiap
        bingkai, \code{FixedUpdate} berlangsung dengan selang tetap untuk
        fisika, dan \code{OnDisable} ketika skrip dimatikan.
  \item \istilah{Scene} adalah satu tingkat permainan yang disimpan sebagai
        berkas. Folder \code{Assets} dan \code{ProjectSettings} disimpan pada
        repositori, sedangkan \code{Library}, \code{Temp}, \code{Obj}, dan
        \code{Build} tidak.
  \item Konsep yang tetap sama persis: kecepatan dinyatakan dalam satuan per
        detik dan dikalikan selang waktu, arah dinormalkan sebelum dikalikan
        kecepatan, posisi disimpan sebagai pecahan, deteksi tabrakan berbeda
        dari penyelesaiannya, dan perhitungan yang mahal sebaiknya disebar.
  \item \code{Time.deltaTime} menggantikan \code{delta_seconds}, dan
        \code{Debug.Log} beserta panel \istilah{Console} menggantikan
        \code{print} beserta terminal.
  \item Pengetahuan tentang cara kerja permainan tidak dapat dipelajari dari
        panduan Unity, sebab panduan itu membahas cara memakai perkakasnya.
        Pengetahuan itulah yang diperoleh dengan menulis seluruh mekanismenya
        sendiri pada blok Python.
\end{itemize}
\end{rangkumanbox}
